\documentclass[11pt]{article}
\usepackage[a4paper,margin=27mm]{geometry}
\usepackage{amsmath,amssymb,amsthm,hyperref}
\newtheorem{theorem}{Theorem}
\newtheorem{lemma}{Lemma}
\newtheorem{corollary}{Corollary}
\newcommand{\Ad}{\operatorname{Ad}}
\newcommand{\ad}{\operatorname{ad}}
\title{Regular fair realizations exist, but require factorially many parameters}
\author{Proof with completed independent mathematical review}
\date{30 September 2026}
\begin{document}\maketitle
Work in the reduced free nilpotent Lie algebra $\mathfrak g_n$ on
$X_1,\ldots,X_n$, where every bracket containing a repeated generator
is zero. Its dimension is
\[
 d_n=\sum_{k=1}^n\binom nk(k-1)!
     \sim e(n-1)!.
\]
A positive word $w=i_1\cdots i_L$ has weighted endpoint map
\[
 \Phi_w(t)=\prod_{a=1}^L\exp(t_aX_{i_a}),\qquad t_a>0.
\]
We call it regular at $t$ when this map has differential rank $d_n$.
This is a statement about varying face weights while keeping their
order fixed, not about changing face labels inside that same order.

\begin{lemma}[A short regular positive prefix]\label{prefix}
There is a positive unit-weight word of length at most $n d_n$ whose
weighted endpoint map is regular at the all-one vector.
\end{lemma}
\begin{proof}
In the reduced algebra $(\ad X_i)^2=0$, so
\[
 \Ad_{\exp X_i}=I+\ad X_i.
\]
Every nested bracket of length at most $n$ is therefore an integer
linear combination of vectors $\Ad_vX_j$, where $v$ is a positive
word of length at most $n-1$. Since nested brackets span
$\mathfrak g_n$, these vectors span the whole Lie algebra.

For a current positive prefix with endpoint $P$, let $S$ be the span
of its existing right-trivialized weight derivatives. A face labeled
$i$ following prefix $Q$ has such derivative $\Ad_QX_i$.
If $S\ne\mathfrak g_n$, the preceding spanning assertion, after
conjugating by $P$, supplies a positive word $v$ of length at most
$n-1$ and a label $j$ with $\Ad_{Pv}X_j\notin S$.
Append $v$ and then one $j$. The newly appended $j$ has that derivative,
so the derivative span increases. Earlier derivatives are unchanged
by appending a suffix. At most $d_n$ steps, each adding at most $n$
letters, give full rank.
\end{proof}

\begin{theorem}[Existence of a regular fair word]\label{fair}
For every $n$, some equal-composition permutation-fair word of total
length $\exp(O(n\log(n+1)))$ is regular for its positive face-weight
endpoint map.
For that word, imposing the $n$ fixed total-weight constraints leaves
rank $d_n-n$ onto the fixed-composition signature manifold.
\end{theorem}
\begin{proof}
Apply the established positive prefix-completion theorem in
\path{order_doubling_prefix_completion.tex} to the
word from Lemma~\ref{prefix}. It extends literally to a finite
permutation-fair equal-composition word. The original prefix derivatives
still have full rank, so the completed word is regular. Its common count
is at most $(n d_n+1)\exp(O(n\log(n+1)))$, and multiplying by $n$
retains total length $\exp(O(n\log(n+1)))$.

The abelianization of the endpoint differential is precisely the
vector of total-weight variations of the individual letters. For a
Lie tangent vector with zero abelianization, any preimage under the
surjective differential already has zero total-weight variation on
each letter. Thus restriction to the fixed-count tangent space is
surjective onto the $d_n-n$ dimensional commutator part.
\end{proof}

\begin{corollary}[A limitation of generic full-rank correction]
Every positive weighted word with full endpoint rank has at least
$d_n=\exp(\Omega(n\log n))$ face-weight parameters, hence at least
that many faces. The same face-count conclusion holds if its fixed-count
differential is surjective onto the full fixed-count signature manifold.
Consequently any hypothetical $\exp(O(n))$-length fair construction
must be singular for this endpoint map. A generic implicit-function
argument that requires surjectivity onto the entire reduced Lie
signature space cannot alone provide such an exponential construction.
\end{corollary}
\begin{proof}
The unrestricted rank is at most the number $L$ of face weights.
With one fixed total per letter, the source dimension is $L-n$ and
surjectivity requires $L-n\ge d_n-n$. Both give $L\ge d_n$.
The dimension asymptotic follows by writing
\[
 \frac{d_n}{(n-1)!}
   =\sum_{j=0}^{n-1}\frac{n}{(n-j)j!}
\]
and splitting the exponential-series tail, or by elementary dominated
estimates on $j\le n/2$ and the factorially small remaining terms.
\end{proof}

\begin{corollary}[Optimal asymptotic order among regular fair words]
Let $L_n^{\mathrm{reg}}$ be the minimum total length of a unit-weight
permutation-fair word whose positive face-weight endpoint map has full
rank. Then
\[
 \log L_n^{\mathrm{reg}}=\Theta(n\log n).
\]
\end{corollary}
\begin{proof}
The dimension bound gives $L_n^{\mathrm{reg}}\ge d_n$, and
Theorem~\ref{fair} supplies the matching logarithmic-order upper bound.
\end{proof}
These results are compatible with singular fair realizations and their
inherited singularity. They determine the asymptotic logarithmic order
only for \emph{regular} fair words; the unrestricted minimum may be
smaller. The regular upper bound uses the order-doubling completion
proved later in this research session.
\end{document}
